% =====================================================================================
% REVISIÓN 1  -  Introducción, Planteamiento del problema, Objetivos, Justificación, Alcances y limitaciones
% Escribe tu texto FUERA de los recuadros \begin{guia} ... \end{guia}.
% No cambies los títulos ni borres los comentarios «% audit:req=...» al final de cada \section / \subsection.
% =====================================================================================

\section{Introducción} \label{sec:introduccion} % audit:req=introduccion

\begin{guia}[Guía: Introducción (Revisión 1)]
\textbf{¿Para qué sirve?} Presenta tu trabajo a un lector que no lo conoce y lo invita a seguir leyendo. Es el resumen
``de entrada'' del proyecto: contexto, problema, propuesta y organización del documento.

\textbf{Debe responder:}
\begin{itemize}[nosep,leftmargin=*]
  \item ¿En qué contexto (sector, empresa, área tecnológica) se desarrolla el proyecto?
  \item ¿Cuál es el problema o la necesidad, dicho en pocas líneas? (el detalle va en \emph{Planteamiento del problema}).
  \item ¿Qué se propone hacer y qué se espera lograr?
  \item ¿Por qué es importante? (una idea; el detalle va en \emph{Justificación}).
  \item ¿Cómo está organizado el documento? (una frase por sección o capítulo).
\end{itemize}

\textbf{Extensión mínima:} 150 palabras (1 a 2 páginas), redactadas con tus propias palabras.

\textbf{Errores frecuentes:} copiar definiciones de Internet; escribir solo ``en este capítulo se verá\ldots'' sin decir
nada concreto; dejar la introducción para el final y no actualizarla cuando el proyecto cambia.

\textbf{Cómo se revisa:} claridad del contexto, que el problema y la propuesta se entiendan sin leer el resto, y que la
organización descrita coincida con las secciones reales del documento.
\end{guia}

% >>> ESCRIBE AQUÍ el texto de la introducción (fuera de la guía) <<<


\subsection{Planteamiento del problema} \label{sec:problema} % audit:req=planteamiento

\begin{guia}[Guía: Planteamiento del problema (Revisión 1)]
\textbf{¿Para qué sirve?} Describe con precisión la situación que quieres resolver y demuestra que es un problema real.

\textbf{Debe responder:}
\begin{itemize}[nosep,leftmargin=*]
  \item ¿Qué situación actual es insatisfactoria o genera una necesidad? Descríbela con hechos, ejemplos o cifras.
  \item ¿Quién la padece (usuarios, área, empresa) y cómo se manifiesta (retrasos, errores, costos, quejas)?
  \item ¿Qué consecuencias tiene no resolverla?
  \item ¿Qué se ha intentado hasta ahora y qué falta (la brecha)? Conecta con los \emph{Antecedentes}.
  \item ¿Cuál es el problema en una sola frase? Escríbelo como enunciado del problema o como pregunta de
        investigación (p. ej. ``¿Cómo reducir X en Y mediante Z?'').
\end{itemize}

\textbf{Extensión mínima:} 120 palabras.

\textbf{Errores frecuentes:} describir la solución en lugar del problema (``el problema es que no hay un sistema web'');
problemas sin evidencia; enunciados tan amplios que no se pueden resolver en una estadía.

\textbf{Cómo se revisa:} que el problema esté delimitado, sea verificable y que los objetivos lo atiendan.
\end{guia}

% >>> ESCRIBE AQUÍ el planteamiento del problema <<<


\subsection{Objetivos} \label{sec:objetivos} % audit:req=objetivos

\begin{guia}[Guía: Objetivo general y objetivos específicos (Revisión 1)]
\textbf{¿Para qué sirve?} Fija qué vas a lograr. Todo lo demás (resultados y conclusiones) se compara contra estos objetivos.

\textbf{Objetivo general} (uno): verbo en infinitivo + qué se hará + con qué o cómo + para qué, y si es posible un criterio
medible. \emph{Ejemplo:} ``Desarrollar un sistema web de notificaciones para familiares de pacientes hospitalizados que
reduzca en 50\% las llamadas de consulta al área de enfermería.''

\textbf{Objetivos específicos} (de 3 a 6): pasos concretos y verificables que, al cumplirse, logran el objetivo general.
Cada uno debe poder demostrarse con un resultado (un módulo, un diagrama, un experimento) que luego aparecerá en
\emph{Resultados} y se retomará en \emph{Conclusiones}. Deben ser SMART: específicos, medibles, alcanzables, relevantes y
con tiempo definido.

\textbf{Debe responder:} ¿qué voy a entregar? ¿cómo sabré que lo logré? ¿cada objetivo específico aporta al general?

\textbf{Extensión mínima:} 40 palabras en total.

\textbf{Errores frecuentes:} verbos no medibles (``conocer'', ``entender'', ``estudiar''); objetivos que son actividades
(``investigar sobre\ldots'') y no logros; más de un objetivo general; objetivos que luego no se atienden en el documento.

\textbf{Cómo se revisa:} coherencia con el problema, medibilidad y que cada objetivo tenga un resultado asociado.
\end{guia}

\subsubsection{Objetivo general} \label{subsec:objetivo_general}
% >>> ESCRIBE AQUÍ el objetivo general (una sola oración) <<<

\subsubsection{Objetivos específicos} \label{subsec:objetivos_especificos}
% >>> ESCRIBE AQUÍ los objetivos específicos como lista numerada: \begin{enumerate} \item ... \end{enumerate} <<<


\subsection{Justificación} \label{sec:justificacion} % audit:req=justificacion

\begin{guia}[Guía: Justificación (Revisión 1)]
\textbf{¿Para qué sirve?} Explica por qué vale la pena dedicar tiempo y recursos a este proyecto.

\textbf{Debe responder:}
\begin{itemize}[nosep,leftmargin=*]
  \item ¿Quién se beneficia y de qué manera? (empresa, usuarios, la comunidad, tú).
  \item ¿Qué se gana en términos prácticos, económicos, sociales, tecnológicos o científicos? Usa datos cuando existan.
  \item ¿Qué pasaría si no se hace?
  \item ¿Por qué es viable hacerlo ahora y por qué con este enfoque?
  \item ¿Cómo se relaciona con el perfil de egreso de tu programa académico?
\end{itemize}

\textbf{Extensión mínima:} 100 palabras.

\textbf{Errores frecuentes:} repetir el problema con otras palabras; afirmaciones sin sustento (``mejorará mucho la
eficiencia''); olvidar al beneficiario.

\textbf{Cómo se revisa:} que los beneficios sean concretos y consistentes con el problema y los objetivos.
\end{guia}

% >>> ESCRIBE AQUÍ la justificación <<<


\subsection{Alcances y limitaciones} \label{sec:alcances_limitaciones} % audit:req=alcances

\begin{guia}[Guía: Alcances y limitaciones (Revisión 1)]
\textbf{¿Para qué sirve?} Marca las fronteras del proyecto para evitar expectativas irreales.

\textbf{Alcances} -- lo que \emph{sí} incluye el trabajo. Debe responder: ¿qué módulos, funciones, procesos, usuarios o
datos se cubren? ¿en qué plataformas o entornos? ¿qué entregables habrá? ¿qué queda \emph{explícitamente fuera}?

\textbf{Limitaciones} -- lo que restringe el trabajo y que no depende totalmente de ti. Debe responder: ¿qué límites hay
de tiempo (un cuatrimestre), recursos, acceso a datos o a la infraestructura de la empresa, tecnología, permisos?
¿cómo pueden afectar la validez o generalización de los resultados?

\textbf{Extensión mínima:} 60 palabras en total.

\textbf{Errores frecuentes:} prometer más de lo que cabe en la estadía; confundir limitaciones con excusas; no decir qué
queda fuera.

\textbf{Cómo se revisa:} realismo del alcance y honestidad de las limitaciones.
\end{guia}

\subsubsection{Alcances} \label{subsec:alcances}
% >>> ESCRIBE AQUÍ los alcances <<<

\subsubsection{Limitaciones} \label{subsec:limitaciones}
% >>> ESCRIBE AQUÍ las limitaciones <<<
